\section{Alcance de los resultados y limitaciones}
\label{sec:disc_limitaciones}
% Extensión estimada: 3 páginas, dos listas y cinco párrafos (ficha 6.5).
% Fuentes: sec:met_eval_generacion, sec:met_estadistica, sec:met_referencias, sec:met_pipeline (piloto de la plantilla),
% sec:met_hiba; sec:res_gen_nivel1, sec:res_gen_ablaciones, sec:res_gen_cualitativos, sec:res_hiba, sec:res_robustez,
% sec:res_sup_clases, sec:res_gen_clases, sec:res_gen_pipeline, sec:res_costo; sec:estado_arte_reportes; secciones
% "Amenazas a la validez" y "Limitaciones" de Docs/experimentacion/experimentacion.tex, con las cifras de la partición
% por lesión. Recibe del Capítulo 1 las limitaciones de alcance. Reemplaza las antiguas 6.4 (métricas) y 6.7.

Las conclusiones de las secciones anteriores valen dentro de las condiciones en que se obtuvieron: imágenes
dermatoscópicas de \emph{HAM10000}, métricas automáticas y el cómputo de una sola GPU. La primera de esas condiciones
que conviene precisar es qué establecen las métricas. La consistencia diagnóstica busca el nombre de la clase en el
texto, de modo que mide la fidelidad diagnóstica del reporte y no si un dermatólogo lo aceptaría, y las métricas con
referencia del Nivel~1 miden sobre todo el parecido de estilo con un texto ajeno (Sección~\ref{sec:res_gen_nivel1}).
En cuanto a la descripción, se evaluó que cambia con la imagen (Sección~\ref{sec:res_gen_ablaciones}) y cuántas
entidades clínicas nombra. Comprobar que cada rasgo descrito esté en la imagen requiere la validación adicional de un
especialista, y los ejemplos de la Sección~\ref{sec:res_gen_pipeline} muestran defectos que estas métricas automáticas
no registran. Esta limitación es compartida por la literatura revisada, en la que la evaluación con
expertos es escasa: cinco dermatólogos en SkinGPT-4~\cite{zhou2024pretrained}, dos entre veinte participantes en
SkinSavvy2~\cite{kim2025skinsavvy2} y un médico en el trabajo de \textcite{chakrabarti2025skin}, mientras que los
demás trabajos evalúan el texto por su parecido con una referencia o por la exactitud del diagnóstico
(Sección~\ref{sec:estado_arte_reportes}).

Además de esa distancia entre lo medido y la validez clínica, conviene tener en cuenta cuatro aspectos de la validez de los resultados:
\begin{itemize}
    \item \textbf{Validez interna.} No todos los experimentos usan el mismo protocolo de inferencia. Los resultados
    principales de clasificación promedian dos vistas de cada imagen con \gls{tta}, mientras que la robustez y la
    interpretabilidad usan una sola vista, y el costo se midió en tres GPU distintas. Cada tabla declara su protocolo,
    pero los valores absolutos de tablas distintas no siempre son comparables. Además, las variantes de la cascada sin
    imagen y con un diagnóstico equivocado usan un subconjunto de 400 imágenes con Mini-VLM y LLaVA-Med
    (Sección~\ref{sec:met_referencias}), de modo que sus intervalos son más amplios que los del conjunto de prueba completo.
    \item \textbf{Validez de la conclusión estadística.} Las comparaciones estadísticas usan como unidad la imagen de prueba y se
    calculan sobre la ejecución de referencia de cada modelo (Sección~\ref{sec:met_estadistica}). Miden, por lo
    tanto, la variación debida a la muestra de imágenes, pero no la variación entre entrenamientos, cuyo rango supera una décima en las clases minoritarias y en el melanoma, y que se informa aparte como desvío entre las diez ejecuciones. La ejecución de referencia, fijada de antemano, queda por encima de la media de las diez, de modo que la cascada se mide con un clasificador algo mejor que el promedio: en melanoma, su sensibilidad es 0.701, frente a 0.585 de media. Las
    clases con pocas imágenes, como \gls{df} con 17, dan intervalos amplios, y algunas diferencias, como la del
    ensamble frente a Swin-Tiny o la de DermLIP con una cabeza frente a los supervisados, no se compararon estadísticamente.
    \item \textbf{Validez de constructo.} La consistencia diagnóstica no juzga el reporte, el Nivel~1 mide sobre todo
    el estilo, y la cobertura de entidades cuenta términos clínicos sin comprobar que correspondan a la imagen. En el
    nivel diferencial, además, la consistencia se parece a acertar entre las dos primeras opciones y no equivale a una
    exactitud.
    \item \textbf{Validez externa.} La evaluación principal usa un solo conjunto, y en \emph{HIBA} Swin-Tiny pierde
    cerca del 40\,\% de su sensibilidad macro, también en las imágenes dermatoscópicas, con un mapeo de diagnósticos
    que agrega ruido en \gls{akiec} (Sección~\ref{sec:res_hiba}). Como la fidelidad diagnóstica del reporte coincide
    con la exactitud del clasificador, es esperable que también caiga en otro conjunto de datos, aunque no se midió. Las
    comparaciones con la literatura, por su parte, ubican el orden de magnitud de los resultados, pero no permiten
    ordenar los modelos.
\end{itemize}

A estos aspectos se suman las limitaciones de alcance del trabajo. La principal es la ausencia de especialistas.
Ningún dermatólogo leyó los reportes, y no existe para las imágenes de \emph{HAM10000} un conjunto de reportes
redactados por especialistas con el que compararlos. Por eso el trabajo mide la fidelidad diagnóstica de los reportes, pero no su validez clínica. Tampoco se pronuncia sobre la fotografía clínica ni sobre la histopatología, donde lo medido no puede trasladarse sin volver a medirlo.

Por último, no todos los resultados tienen el mismo respaldo. Las decisiones de selección se tomaron en validación: el modelo, por la media de las ejecuciones, y el umbral de melanoma, con una especificidad mínima fijada de antemano. El conjunto de prueba de la partición por lesión y \emph{HIBA} aportan, entonces, una evaluación independiente de esas decisiones, y el ruteo no ajusta ningún umbral sobre el conjunto de prueba. La única excepción es la instrucción \emph{noreasoning}, que se eligió en un piloto sobre 100 imágenes, de las cuales 17 pertenecen al conjunto de prueba (Sección~\ref{sec:met_pipeline}), sin fijar ningún umbral sobre él. En cambio, las explicaciones que este capítulo propone, como la jerarquía entre las fuentes del reporte o el origen de los errores compartidos, y los análisis que se hicieron después de ver los
resultados, como la clasificación de las fallas de los reportes o los ejemplos, describen el conjunto de prueba sin haberse
confirmado en otros datos, y quedan como exploratorios. Con estas limitaciones delimitadas, el
Capítulo~\ref{cp:conclusiones} responde las preguntas de investigación y reúne las conclusiones del trabajo.
